\section{Methods}

\subsection{Study design and analysis plan}
Norman et al.\ 2019 \citep{norman2019} served as the discovery dataset: all modelling and evaluation choices were developed on it, and several baselines and metrics were added after first results on its folds (disclosed in the Supplement). Before any other dataset was preprocessed or analysed, an analysis plan was written that fixed the datasets, inclusion rules, preprocessing, splits, models, hyperparameter grids, confirmatory tests and multiple-testing correction (\texttt{docs/analysis\_plan.md}; full SHA-256 in Supplementary Table~\ref{tab:hashes}). One addendum, adding a metric-calibration analysis motivated by \citet{miller2026}, was written before any confirmation dataset was analysed (\texttt{docs/analysis\_plan\_addendum.md}). Both hashes are verified and recorded by the study runner. The plan and addendum were not registered with a third party, so their timing rests on the project's own records.

\subsection{Data}
All data were the harmonised H5AD files of scPerturb \citep{peidli2024}, Zenodo record 13350497 (version 1.4, CC-BY-4.0), each verified against its published MD5 checksum. Candidate confirmation datasets were Replogle et al.\ 2022 K562 essential-gene and RPE1 CRISPRi screens \citep{replogle2022}, the Adamson et al.\ 2016 unfolded-protein-response CRISPRi screen \citep{adamson2016}, the Dixit et al.\ 2016 day-7 K562 knockout screen \citep{dixit2016}, and the Tian et al.\ 2021 iPSC-derived neuron CRISPRa and CRISPRi screens \citep{tian2021}.

\paragraph{Inclusion.} A dataset was included if it passed checksum and schema validation; had at least 100 cells labelled \texttt{control} in the harmonised \texttt{perturbation} column; had at least 30 single-gene conditions with at least 50 cells (25 if fewer than 30 qualified at 50); and at least 80\% of its target labels resolved to a measured gene, directly or via HGNC-to-Ensembl mapping. Control groups were never inferred from guide names. Conditions with more than one target were excluded from the confirmation study.

\subsection{Preprocessing and response targets}
Where present, cells were restricted to \texttt{good\_coverage} and \texttt{number\_of\_cells}$=1$. Counts were scaled to $10^4$ per cell over all genes and log1p-transformed. Statistics were accumulated in chunks along the matrix's storage axis, so peak memory stayed below 2\,GB for the largest dataset (310{,}385 cells). Genes were selected using control cells only. Among genes detected in at least 1\% of control cells, dispersion was computed as the variance divided by the mean of log-normalised expression. Genes were divided into 20 quantile bins of mean expression, dispersion was $z$-scored within each bin, and the 2{,}000 genes with the highest $z$-scores were kept. All target genes were then added.

For condition $p$, gene $g$ and batch $r$ (10x lane), with $w_{pr}$ the fraction of $p$'s cells in batch $r$, the response is
\begin{equation}
\Delta_{pg} \;=\; \sum_r w_{pr}\big(\bar x_{pgr} - \bar c_{gr}\big) \;=\; \bar x_{pg} - \sum_r w_{pr}\,\bar c_{gr},
\end{equation}
where $\bar x$ and $\bar c$ are mean log-normalised expression in perturbed and control cells. The second form avoids storing per-condition-per-batch statistics. Control cells were split at random into two halves within each batch. Training targets $\Delta^{\mathrm{fit}}$ use one half and evaluation targets $\Delta^{\mathrm{eval}}$ the other, so no control noise is shared between the two. The perturbed gene's own entry is masked in its own condition, both in training and in all metrics, because CRISPRa and CRISPRi alter their targets directly.

\subsection{Splits}
In the \emph{unseen-gene} task, target genes are partitioned into five folds. Every condition containing a held-out gene is removed from training and tested. A calibration set is built the same way from 15\% of the remaining genes and used for hyperparameter selection, early stopping and interval calibration. An automated check rejects any fold in which a held-out or calibration gene, or a duplicated gene set, appears in training. In Norman, the \emph{unseen-pair} task holds out pairs whose two single-gene responses are in training.

\subsection{Models}
\textbf{Mean}: the mean training response (per number of targets). \textbf{Ridge}: ridge regression on the residual from the mean, using a perturbation representation. \textbf{Gene-loading ridge}: following \citet{ahlmanneltze2025}, the target is represented by its loadings in a PCA of the training response matrix. \textbf{kNN}: average residual of the nearest training perturbations. \textbf{MLP ensemble}: five bootstrap-trained networks (two hidden layers, 128 and 64 units, dropout 0.1, Huber loss, AdamW, early stopping), predicting 20 principal components of training residuals. \textbf{Additive}, \textbf{scaled additive} (one global factor fitted on training pairs) and \textbf{additive plus mean residual} were used for pairs, together with ridge and MLP variants that learn the residual from the additive prediction.

A perturbation representation is the sum of its targets' gene vectors, plus the targets' control expression and detection rate. Gene vectors were of three kinds. The first was a 32-dimensional co-expression embedding. Log-normalised expression of the selected genes in at most 8{,}000 randomly sampled control cells was centred and scaled to unit variance per gene; genes with zero variance were set to 0. The resulting matrix was decomposed by randomised truncated SVD (seed 0), and each gene's vector is its row of the right singular vectors, multiplied by the singular values and divided by $\sqrt{n_{\text{cells}}}$. The target's own baseline features are its mean log-normalised expression and its detection rate (fraction of cells with nonzero counts) over all control cells. These features use control cells from both random halves, including the half behind the evaluation targets; a post-hoc check recomputed them from the training half alone (Supplement). The second kind of gene vector was the same embedding with rows randomly permuted among genes, as a negative control. Permutation leaves each target's baseline features unchanged, so this control isolates the contribution of co-expression. The third kind was frozen 200-dimensional Gene2vec embeddings \citep{du2019gene2vec}. Perturbation features were standardised using training conditions only within each fold. Full hyperparameter grids and seeds are in Supplementary Table~\ref{tab:grids}. In the discovery dataset, STRING-derived features \citep{szklarczyk2023} were also examined (Supplement).

\subsection{Metrics}
For each held-out perturbation, excluding masked entries: mean squared error (MSE); \emph{skill} $=1-\mathrm{SSE}_{\text{model}}/\mathrm{SSE}_{\text{mean}}$; Pearson correlation of the raw response; and the Pearson and Spearman correlation of the \emph{perturbation-specific} response, the deviation from the training mean \citep[cf.\ the perturbed-centroid reference of][]{vinastorne2025}. The mean baseline has skill 0 by definition and no defined specific correlation. Perturbations are the unit of inference. Intervals are 95\% percentile bootstraps over perturbations (2{,}000 resamples), and paired comparisons use per-perturbation differences.

\paragraph{Technical reproducibility limits.} Each response was re-estimated from two disjoint sets of batches (alternating lanes), each against its own batch-matched controls. The correlation $r$ between the two estimates of the perturbation-specific response gives a Spearman--Brown \citep{spearman1910,brown1910} reliability $2r/(1+r)$ for the full estimate. Its square root approximates the correlation that a noiseless prediction of the expected response \emph{in this experiment} could reach. This rests on three assumptions: the two halves estimate the same expected response with independent errors; the errors have equal variance; and averaging the halves doubles the effective sample. Because the batches are technical lanes within one experiment, the limit captures sampling and lane-to-lane technical noise, but not variation between independent biological replicates, cultures or days. We therefore call it a \emph{technical} reproducibility limit; against an independently repeated experiment, the attainable correlation would be lower.

\paragraph{Metric calibration.} Following the logic of \citet{miller2026}, each metric was scored for a positive control (a technical duplicate: the response from one random half of a condition's cells and controls, predicting the other half) and a negative control (the mean of other conditions). A metric was called calibrated on a dataset if the duplicate won for at least 80\% of conditions.

\subsection{Confirmatory tests}
\textbf{C1:} For each confirmation dataset and pre-specified learned model (ridge, gene-loading ridge, MLP ensemble; co-expression features), a one-sided Wilcoxon signed-rank test of lower per-perturbation MSE than the mean baseline, with Holm correction \citep{holm1979} across all dataset--model tests. A model beats the mean if its adjusted $p<0.05$ \emph{and} the bootstrap interval of the mean MSE difference lies below 0. The two parts test different things. The signed-rank test concerns the location of the per-perturbation differences (whether the model is typically better), and the bootstrap interval concerns the average effect (whether expected error is lower). Requiring both guards against declaring an improvement when most perturbations improve slightly but a few deteriorate enough to raise average error, which matters for a tool whose errors are paid per experiment. \textbf{C2:} co-expression versus permuted co-expression; informative if the 95\% interval of the paired skill difference lies above 0. The plan applied no multiplicity correction to C2. A Holm-corrected two-sided Wilcoxon version is reported as a post-hoc sensitivity analysis. \textbf{C3:} Gene2vec versus co-expression, two-sided Wilcoxon with Holm correction.

\subsection{Exploratory analyses}
\textbf{Training-set size} (Replogle screens): training conditions were subsampled to 25--800 (three draws per size) within fixed folds. \textbf{Cross-context transfer}: for targets measured in both Replogle screens, on shared genes, one cell line's measured response predicted the other's, either raw or as the target dataset's mean plus a rescaled deviation (scale fitted on other targets). \textbf{Pair coefficients} (Norman): for each pair, $\beta$ in $\Delta_{\text{pair}}\approx\beta(\Delta_A+\Delta_B)$, estimated by least squares and by an instrumental-variable estimator, $\hat\beta_{\text{IV}}=\langle s^{(1)},y\rangle/\langle s^{(1)},s^{(2)}\rangle$. Here $s^{(1)},s^{(2)}$ are the summed single responses from disjoint halves of the single-perturbation cells, which removes attenuation from measurement error in the singles.

\subsection{Post-hoc analyses added in revision}
All of the following were added after the confirmatory results were seen. The first was a feature ablation, which refit ridge and the MLP on the primary folds with co-expression only, baseline mean only, detection rate only, both baseline features, and the primary features recomputed from training-half control cells. The second was a magnitude-stratified interval calibration, in which the interval quantile is calibrated separately within five strata of predicted deviation from the training mean, ranked within each perturbation. The strata are therefore defined from predictions alone and are available at prediction time. The third analysed the distribution of per-perturbation MSE differences behind C1.

\subsection{Software and reproducibility}
The pipeline is a Python package (Python 3.11.15, PyTorch 2.14.1, scikit-learn 1.9.1, NumPy 2.4.6, SciPy 1.17.1, pandas 2.3.3, AnnData 0.12.19; exact versions in \texttt{requirements-lock.txt}) with 149 automated tests, run on a laptop CPU (Apple M3, 8\,GB). Two complete runs of the discovery benchmark gave identical per-perturbation metrics. Every run records its configuration, split assignments, data checksums, target digests, package versions and the analysis-plan hashes.
